\chapter{Computational Appendix: Reproducibility and Data Transparency}

\section{Appendix A: Unsupervised Phonetic Discovery Model}
This appendix provides the Python implementation of the Principal Component Analysis (PCA) used to identify emergent latent structures in the phonological dataset. 

\begin{verbatim}
import pandas as pd
import numpy as np
from sklearn.decomposition import PCA

def run_real_unsupervised_discovery():
    # Phonetic Feature Matrix: [Voice, SpreadGlottis, ConstrictedGlottis, 
    # Labial, Coronal, Dorsal]
    
    # Aligned sets: [Vedic, Tamil, Tibetan] (18 dimensions total)
    set_P = [1,0,0, 1,0,0,  0,0,0, 1,0,0,  0,1,0, 1,0,0] # Labials
    set_T = [1,0,0, 0,1,0,  0,0,0, 0,1,0,  0,1,0, 0,1,0] # Dentals
    set_K = [1,0,0, 0,0,1,  0,0,0, 0,0,1,  0,1,0, 0,0,1] # Velars
    set_M = [1,0,0, 1,0,0,  1,0,0, 1,0,0,  1,0,0, 1,0,0] # Nasals (Control)
    set_S = [0,0,0, 0,1,0,  0,0,0, 0,1,0,  0,0,0, 0,1,0] # Fricatives (Control)
    
    X = np.array([set_P, set_T, set_K, set_S, set_M])
    
    pca = PCA(n_components=2)
    latent_space = pca.fit_transform(X)
    
    explained_var = pca.explained_variance_ratio_
    print(f"Explained Variance (PC1): {explained_var[0]*100:.2f}%")
    
if __name__ == "__main__":
    run_real_unsupervised_discovery()
\end{verbatim}

\section{Appendix B: Zero-Shot Predictive Validation Model}
The following script documents the Hidden Markov Model (HMM) configuration used for the predictive validation on the Mahadevan/Wells corpus.

\begin{verbatim}
from hmmlearn import hmm
from sklearn.model_selection import train_test_split
import numpy as np

# model hyperparameters
n_states = 4
model = hmm.CategoricalHMM(n_components=n_states, n_iter=200, random_state=42)

# Training protocol:
# 1. Purge all semantic labels.
# 2. Split 80/20 train-test.
# 3. Apply Laplacian Smoothing (epsilon = 1e-6)
# 4. Measure Log-Likelihood on unseen test set.
\end{verbatim}

\section{Appendix C: Representative Archaeological Data Sample}
A sample of the cleaned structural corpus utilized for the HMM analysis. Data is presented as sequences of numeric sign identifiers (M-codes).

\begin{table}[h]
\centering
\begin{tabular}{ll} \toprule
\textbf{Seal ID} & \textbf{Structural Sequence (Sign IDs)} \\ \midrule
1.1 & 410, 017 \\
3.1 & 405, 017 \\
5.1 & 740, 235 \\
6.1 & 740, 390, 590 \\
7.1 & 368, 390, 125, 033 \\
8.1 & 740, 904, 033, 705, 235 \\
11.1 & 072, 170, 740, 840, 013 \\ \bottomrule
\end{tabular}
\caption{Sample Inscriptions from the Cleaned IVS Corpus.}
\end{table}
